% TOP_PROBLEM: {"id":30007612,"problem_number":"LOCAL-30007612","title":"Exchange-rate forecasting","category_id":21,"category":{"id":21,"name":"economics","display_name":"Economics","description":"Open economic theory statements, published research agendas, and proposed scoped specifications; question provenance and historical priority are qualified per record.","slug":"economics","order_index":21,"created_at":"2026-10-08T00:00:00Z"},"status":"research_question","external_url":null,"legacy_ids":[],"published":false,"statement":"In the explicitly specified setting: Which exchange-rate forecasting methods reliably outperform a random walk out of sample across currencies and horizons? The mathematical statement above fixes the target, assumptions and scope of this catalogue formulation.","economics":{"original_id":101,"field":"International finance and exchange rates","kind":"empirical","formalization_basis":"adapted_research_question","source_status":"explicit_open","priority_status":"earliest_found","priority_note":"A June 1981 IFDP version predates the verified journal abstract. The 1983 finding is model-, horizon- and sample-specific; priority for the generalization question is not established.","earliest_verified_reference":{"authors":["Richard A. Meese","Kenneth Rogoff"],"title":"Empirical Exchange Rate Models of the Seventies: Do They Fit Out of Sample?","year":1983,"url":"https://www.sciencedirect.com/science/article/pii/002219968390017X","doi":"10.1016/0022-1996(83)90017-X","locator":"Abstract, journal p.3","evidence_summary":"Frames the out-of-sample random-walk comparison for specified models."},"current_reference":{"authors":["Bas B. Bakker"],"title":"Reconciling Random Walks and Predictability: A Dual-Component Model of Exchange Rate Dynamics","year":2024,"url":"https://www.imf.org/-/media/files/publications/wp/2024/english/wpiea2024252-print-pdf.pdf","doi":"10.5089/9798400295034.001","locator":"Abstract (PDF p.2); Section 9, pp.35–38; Section 10, p.39","evidence_summary":"Reports forecast improvements in particular countries and horizons."},"formalization_note":"The historical random-walk result is a resolved comparison within a specified sample, not an impossibility theorem. Current papers improve forecasts in subcases; robust transportability is the adapted task."},"statusQualification":"Published question or research agenda verified at the cited locator. The mathematical specification is adapted for this catalogue and includes additional assumptions; source publication does not establish first-ever priority or certify this exact model remains unsolved. The historical random-walk result is a resolved comparison within a specified sample, not an impossibility theorem. Current papers improve forecasts in subcases; robust transportability is the adapted task.","statusReviewedAt":"2026-10-08"}
% FIRST_POSTED: 2026-10-08
% ENTRY_KIND: statement_only
% !TeX program = lualatex
\documentclass[11pt]{article}
\usepackage{fontspec}
\usepackage[a4paper,margin=25mm]{geometry}
\usepackage{amsmath,amssymb,mathtools}
\usepackage{xurl}
\usepackage[hidelinks]{hyperref}
\newcommand{\sref}[2]{\hyperlink{source:#1:#2}{[S#2]}}
\setlength{\emergencystretch}{3em}
\begin{document}
% problemId:30007612
\section*{Exchange-rate forecasting}\label{30007612}
\subsection{Definitions and mathematical statement}
Fix a prespecified set of currency pairs \(C\), forecast horizons \(H\), evaluation dates \(T\), and an information set available in real time. Let \(s_{ct}\) be the logarithm of currency pair \(c\)'s spot exchange rate. A forecasting procedure \(f\) maps dated observations to \(\widehat s^f_{c,t+h}\), with parameter estimation using only information available by \(t\). The no-drift random-walk benchmark is \(\widehat s^{RW}_{c,t+h}=s_{ct}\). For a declared loss function \(\ell\), define
\[\Delta_{ch}(f)=E[\ell(s_{c,t+h},\widehat s^f_{c,t+h})-\ell(s_{c,t+h},s_{ct})].\]
The task is to identify methods and economic conditions with negative, statistically supported loss differences on untouched observations. Specify whether the claim requires improvement for every \((c,h)\in C\times H\) or only a weighted aggregate; these are different research targets. Account for predictor revisions, model selection, overlapping forecast errors and structural breaks. Compare statistical accuracy with feasible trading gains after costs rather than treating them as identical. Derive confidence sets or predictive guarantees under stated stability assumptions and test sensitivity across regimes. A method that beats the benchmark in one sample is a resolved subcase, not a universal solution; neither the classic comparison nor the present task asserts impossibility of improved forecasting.

\par\textbf{Formulation and scope.} The historical random-walk result is a resolved comparison within a specified sample, not an impossibility theorem. Current papers improve forecasts in subcases; robust transportability is the adapted task.

\par\textbf{Open-question provenance.} An explicit question or research agenda was verified at the source locator. This does not by itself certify that every later version remains unresolved \sref{30007612}{1}.

\par\textbf{Historical priority.} A June 1981 IFDP version predates the verified journal abstract. The 1983 finding is model-, horizon- and sample-specific; priority for the generalization question is not established.

\subsection{Short English statement}
In the explicitly specified setting: Which exchange-rate forecasting methods reliably outperform a random walk out of sample across currencies and horizons? The mathematical statement above fixes the target, assumptions and scope of this catalogue formulation.

\subsection{Sources}
\begin{itemize}\raggedright
\item[\textnormal{[S1]}]\hypertarget{source:30007612:1}{} Richard A. Meese, Kenneth Rogoff. 1983. Empirical Exchange Rate Models of the Seventies: Do They Fit Out of Sample?. Abstract, journal p.3.
\url{https://www.sciencedirect.com/science/article/pii/002219968390017X}.
DOI: 10.1016/0022-1996(83)90017-X.
{\small\textit{Earliest verified question/agenda reference; historical priority qualified below. Frames the out-of-sample random-walk comparison for specified models.}}

\item[\textnormal{[S2]}]\hypertarget{source:30007612:2}{} Bas B. Bakker. 2024. Reconciling Random Walks and Predictability: A Dual-Component Model of Exchange Rate Dynamics. Abstract (PDF p.2); Section 9, pp.35–38; Section 10, p.39.
\url{https://www.imf.org/-/media/files/publications/wp/2024/english/wpiea2024252-print-pdf.pdf}.
DOI: 10.5089/9798400295034.001.
{\small\textit{Supporting or current literature; see evidence qualification. Reports forecast improvements in particular countries and horizons.}}
\end{itemize}
\end{document}
